\documentclass[10pt]{article}
% Surcouche compacte pour les interros
% Socle commun dérivé du framework des khôlles SI
% Version autonome et compatible pdfLaTeX pour le dépôt Interros.
\RequirePackage[T1]{fontenc}
\RequirePackage[utf8]{inputenc}
\RequirePackage[french]{babel}
\RequirePackage{lmodern}
\RequirePackage[margin=1.8cm]{geometry}
\RequirePackage{xcolor}
\RequirePackage{graphicx}
\RequirePackage{amsmath,amssymb}
\RequirePackage{tikz}
\RequirePackage{tabularx}
\RequirePackage{array}
\RequirePackage{enumitem}
\RequirePackage[most]{tcolorbox}
\RequirePackage{fancyhdr}
\RequirePackage{hyperref}
\RequirePackage{titlesec}
\usetikzlibrary{arrows.meta,positioning,calc,fit,chains,decorations.pathmorphing,decorations.markings,shapes,shapes.multipart,patterns,graphs,quotes}
\definecolor{SIblue}{RGB}{31,78,121}
\definecolor{SIlightblue}{RGB}{229,240,250}
\definecolor{SIred}{RGB}{190,0,0}
\definecolor{SIgray}{RGB}{235,235,235}
\hypersetup{colorlinks=true,linkcolor=SIblue,urlcolor=SIblue}
\pagestyle{fancy}
\setlength{\headheight}{14pt}
\fancyhf{}
\lhead{Sciences industrielles pour l'ingénieur}
\rhead{Interros ATS}
\cfoot{\thepage}
\titleformat{\section}{\Large\bfseries\color{SIblue}}{}{0pt}{}
\titleformat{\subsection}{\large\bfseries\color{SIblue}}{}{0pt}{}
\setlist[itemize]{leftmargin=*,itemsep=2pt,topsep=2pt}
\setlist[enumerate]{leftmargin=*,itemsep=4pt,topsep=4pt}
\newtcolorbox{donnees}[1][]{enhanced,breakable,colback=white,colframe=SIblue,title=Données,fonttitle=\bfseries,#1}
\newtcolorbox{attention}[1][]{enhanced,breakable,colback=orange!8,colframe=orange!80!black,title=Attention,fonttitle=\bfseries,#1}
\newtcolorbox{methode}[1][]{enhanced,breakable,colback=green!5,colframe=green!45!black,title=Méthode,fonttitle=\bfseries,#1}
\newcommand{\torseur}[2]{\left\{\mathcal V_{#1}\right\}_{#2}}
% Bibliothèques SI simplifiées, compatibles avec les interros
\providecommand{\og}{``}
\providecommand{\fg}{''}

% Graphe de liaisons
\tikzset{
  glsolide/.style={circle,draw=SIblue,very thick,minimum size=10mm,fill=SIlightblue,font=\bfseries},
  glliaison/.style={draw=SIred,very thick,rounded corners=2pt,font=\small,align=center}
}
\newenvironment{grapheliaisons}[1][]{\begin{center}\begin{tikzpicture}[node distance=25mm,#1]}{\end{tikzpicture}\end{center}}
\newcommand{\GLsolide}[3]{\node[glsolide] (#1) at (#3) {#2};}
\newcommand{\GLliaison}[4][]{\draw[glliaison,#1] (#2) -- node[midway,fill=white,inner sep=2pt]{#4} (#3);}
\newcommand{\GLpivot}[4][]{\GLliaison[#1]{#2}{#3}{Pivot\\$#4$}}
\newcommand{\GLglissiere}[4][]{\GLliaison[#1]{#2}{#3}{Glissière\\$#4$}}
\newcommand{\GLencastrement}[3][]{\GLliaison[#1]{#2}{#3}{Encastrement}}
\newcommand{\GLhelicoidale}[4][]{\GLliaison[#1]{#2}{#3}{Hélicoïdale\\$#4$}}

% Schémas cinématiques simplifiés
\newenvironment{schemacinematique}[1][]{\begin{center}\begin{tikzpicture}[line cap=round,line join=round,#1]}{\end{tikzpicture}\end{center}}
\newcommand{\SCbati}[2]{\coordinate (#1) at (#2);\draw[thick] ($(#1)+(-.65,0)$)--($(#1)+(.65,0)$);\foreach \x in {-.55,-.3,...,.55}{\draw ($(#1)+(\x,0)$)--++(-.18,-.18);}}
\newcommand{\SCpivot}[2]{\node[draw,circle,minimum size=5mm,thick] (#1) at (#2) {};}
\newcommand{\SCglissiere}[2]{\node[draw,rectangle,minimum width=10mm,minimum height=5mm,thick] (#1) at (#2) {};\draw[thick] ($(#1)+(-.35,-.18)$)--($(#1)+(.35,.18)$);}
\newcommand{\SChelicoidale}[2]{\node[draw,circle,minimum size=5mm,thick] (#1) at (#2) {};\draw[thick] ($(#1)+(-.15,0)$) arc (180:-90:.15);}
\newcommand{\SCrelier}[3][]{\draw[very thick,#1] (#2)--(#3);}

% SysML simplifié
\newenvironment{diagrammeSysML}[1][]{\begin{center}\begin{tikzpicture}[node distance=18mm,#1]}{\end{tikzpicture}\end{center}}
\newcommand{\SYSacteur}[4][]{\node[draw,circle,minimum size=3mm,#1] (head#4) at (#2) {};\draw (head#4)--++(0,-.55) coordinate (body#4);\draw ($(body#4)+(0,.35)$)--++(-.28,-.22);\draw ($(body#4)+(0,.35)$)--++(.28,-.22);\draw (body#4)--++(-.25,-.35);\draw (body#4)--++(.25,-.35);\node[below=8mm of head#4] (#4) {#3};}
\newcommand{\SYScas}[4][]{\node[draw,ellipse,align=center,minimum width=28mm,minimum height=10mm,#1] (#4) at (#2) {#3};}
\newcommand{\SYSinclude}[2]{\draw[dashed,->] (#1)--node[above]{\scriptsize\og include\fg}(#2);}
\newcommand{\SYSextend}[2]{\draw[dashed,->] (#1)--node[above]{\scriptsize\og extend\fg}(#2);}

% Tableau de mobilités
\newcommand{\mobilites}[6]{\begin{tabular}{c|c}R&T\\\hline #1&#2\\#3&#4\\#5&#6\end{tabular}}


\geometry{a4paper,left=13mm,right=13mm,top=18mm,bottom=16mm,headheight=28pt,headsep=8mm,footskip=10mm}
\fancyhf{}
\renewcommand{\headrulewidth}{0pt}
\renewcommand{\footrulewidth}{0.6pt}
\fancyfoot[L]{\footnotesize Lycée Robert Doisneau}
\fancyfoot[R]{\footnotesize \thepage}
\newcommand{\InterroHeader}[2]{%
\noindent\begin{tcolorbox}[enhanced,boxrule=.7pt,arc=2mm,colback=SIblue,colframe=SIblue,left=3mm,right=3mm,top=1.5mm,bottom=1.5mm]
\begin{minipage}[c]{.18\linewidth}\color{white}\bfseries\Large SI\end{minipage}%
\begin{minipage}[c]{.60\linewidth}\centering\color{white}\bfseries\large #1\end{minipage}%
\begin{minipage}[c]{.20\linewidth}\raggedleft\color{white}\bfseries Interro\\de cours\end{minipage}
\end{tcolorbox}
\vspace{-1mm}
\noindent\begin{minipage}[c]{.78\linewidth}\begin{tcolorbox}[height=12mm,colback=white,colframe=black,boxrule=1pt,arc=1.5mm,left=2mm,top=2mm]\textbf{Nom :}\end{tcolorbox}\end{minipage}\hfill
\begin{minipage}[c]{.13\linewidth}\begin{tcolorbox}[height=12mm,colback=white,colframe=black,boxrule=1pt,arc=2mm,halign=center,valign=center]\textbf{#2}\end{tcolorbox}\end{minipage}
\vspace{2mm}}
\newcounter{interroquestion}
\newcommand{\question}[2]{\refstepcounter{interroquestion}\par\medskip\noindent\textbf{Question \theinterroquestion :}\quad #1\hfill\textbf{(#2 points)}\par\smallskip}
\newtcolorbox{reponse}[1][35mm]{colback=white,colframe=black,boxrule=.5pt,arc=0mm,height=#1}
\newcommand{\resetquestions}{\setcounter{interroquestion}{0}}
\newcommand{\chaineenergie}{\begin{tikzpicture}[font=\small,>=Latex,node distance=4mm,bloc/.style={draw,minimum width=25mm,minimum height=10mm,align=center}]
\node[bloc](a){\dots};\node[bloc,right=of a](b){MODULER};\node[bloc,right=of b](c){\dots};\node[bloc,right=of c](d){\dots};\node[bloc,fill=yellow!30,right=of d](e){AGIR};
\draw[->,thick](a)--(b);\draw[->,thick](b)--(c);\draw[->,thick](c)--(d);\draw[->,thick](d)--(e);\end{tikzpicture}}
\newcommand{\chaineinformation}{\begin{tikzpicture}[font=\small,>=Latex,node distance=5mm,bloc/.style={draw,minimum width=25mm,minimum height=11mm,fill=yellow!30,align=center}]
\node[bloc](a){\dots};\node[bloc,right=of a](b){\dots};\node[bloc,right=of b](c){\dots};\draw[->,very thick,SIblue](a)--(b);\draw[->,very thick,SIblue](b)--(c);\end{tikzpicture}}

\setlength{\parindent}{0pt}
\usetikzlibrary{circuits.ee.IEC,positioning,arrows.meta}

\tikzset{every circuit symbol/.style={draw,thick}, circuit declare symbol=resistance,
  set resistance graphic={draw,minimum width=7mm,minimum height=2.6mm}}

\newcommand{\circuitcondensateur}{%
\begin{tikzpicture}[>=Latex,scale=.9]
  \draw (0,0)--(0,0.65);\draw (-.38,.65)--(.38,.65);\draw (-.38,.9)--(.38,.9);\draw (0,.9)--(0,1.6);
  \draw[->,red,thick] (-.8,1.55)--(-.8,.95) node[midway,left] {$i(t)$};
  \draw[->,blue!70!black,thick] (.8,.15)--(.8,1.5) node[midway,right] {$u(t)$};
  \node[left] at (-.05,.78) {$C$};
\end{tikzpicture}}

\newcommand{\circuitqfive}{%
\begin{tikzpicture}[scale=.72]
  \draw (0,0)--(6,0) node[right] {$B$}; \draw (0,3)--(6,3) node[right] {$A$};
  \foreach \x/\r/\e in {1/R_1/E_1,3/R_2/E_2}{
    \draw (\x,3)--(\x,2.4);\draw (\x-.18,2.4) rectangle (\x+.18,1.6);\node[right] at (\x+.2,2) {$\r$};
    \draw (\x,1.6)--(\x,.85);\draw (\x-.28,.85)--(\x+.28,.85);\draw (\x-.17,.62)--(\x+.17,.62);\draw (\x,.62)--(\x,0);\node[right] at (\x+.25,.72) {$\e$};}
  \draw (5,3)--(5,2.4);\draw (4.82,2.4) rectangle (5.18,1.6);\node[right] at (5.2,2) {$R_3$};\draw (5,1.6)--(5,0);
\end{tikzpicture}}

\newcommand{\circuitqsix}{%
\begin{tikzpicture}[scale=.72]
  \node[left] at (0,2) {$A$};\node[left] at (0,0) {$B$};
  \draw (0,2)--(.7,2);\draw (.7,1.82) rectangle (1.8,2.18);\node[above] at (1.25,2.2) {$R$};\draw (1.8,2)--(4.8,2);\draw (0,0)--(4.8,0);
  \foreach \x in {2.2,3.3,4.4}{\draw (\x,2)--(\x,1.55);\draw (\x-.18,1.55) rectangle (\x+.18,.45);\node[left] at (\x-.2,1) {$R$};\draw (\x,.45)--(\x,0);}
\end{tikzpicture}}

\newcommand{\circuitqseven}{%
\begin{tikzpicture}[scale=.68,>=Latex]
  \draw (0,0)--(5,0);\draw (0,0)--(0,.55);\draw (0,2.45)--(0,3)--(.8,3);
  \draw (0,.55) circle (.42);\draw[->] (0,.1)--(0,2.85) node[midway,left] {$E$};
  \draw (.8,2.82) rectangle (2,3.18);\node[above] at (1.4,3.18) {$R_1$};\draw (2,3)--(5,3);
  \draw (2.5,3)--(2.5,2.55);\draw (2.32,2.55) rectangle (2.68,.45);\node[left] at (2.28,1.5) {$R_2$};\draw (2.5,.45)--(2.5,0);
  \draw (3.35,2.82) rectangle (4.2,3.18);\node[above] at (3.78,3.18) {$R$};
  \draw (5,3)--(5,2.55);\draw (4.82,2.55) rectangle (5.18,.45);\node[right] at (5.2,1.5) {$R$};\draw (5,.45)--(5,0);
  \draw[->] (2.9,.55)--(2.9,2.45) node[midway,right] {$U_2$};
  \draw[->] (1.9,2.55)--(.9,2.55) node[midway,below] {$U_1$};
\end{tikzpicture}}

\newcommand{\circuitqeight}{%
\begin{tikzpicture}[scale=.72,>=Latex]
  \draw (0,0)--(5,0);\draw (0,3)--(5,3);
  \draw (0,0)--(0,.85);\draw (0,2.15)--(0,3);\draw (0,1.5) circle (.55);\draw (-.35,1.5)--(.35,1.5);
  \draw[->,thick] (-.8,.4)--(-.8,2.6) node[above] {$i=1{,}4\,\mathrm{A}$};
  \foreach \x/\r in {1.5/R_1,3/R_2,4.5/R_3}{\draw (\x,3)--(\x,2.5);\draw (\x-.22,2.5) rectangle (\x+.22,.5);\node at (\x,1.5) {$\r$};\draw (\x,.5)--(\x,0);}
  \draw[->,thick] (1.5,3.2)--(1.5,2.7) node[right] {$i_1$};
\end{tikzpicture}}

\newcommand{\spectrequest}{%
\begin{tikzpicture}[x=.0015cm,y=.35cm,>=Latex]
  \draw[->] (0,0)--(5400,0) node[right] {$f\,(\mathrm{Hz})$};\draw[->] (0,0)--(0,14) node[above] {$G\,(\mathrm{V})$};
  \foreach \x in {0,1000,...,5000}{\draw (\x,.15)--(\x,-.15) node[below] {\scriptsize\x};}
  \foreach \y in {0,2,...,14}{\draw (60,\y)--(-60,\y) node[left] {\scriptsize\y};\draw[gray!40] (0,\y)--(5200,\y);}
  \foreach \x/\y in {500/12.5,1500/4,2400/1.2,3400/3,4400/3.8}{\draw[blue!70!black,very thick] (\x,0)--(\x,\y);}
\end{tikzpicture}}

\begin{document}
\InterroHeader{Outils et lois de l'électrocinétique}{20}

\question{Donner au moins cinq critères permettant de choisir un capteur.}{1}
\begin{reponse}[20mm]\end{reponse}
\question{Donner des solutions technologiques de capteurs permettant de mesurer une masse et une position angulaire.}{1}
\begin{reponse}[20mm]\end{reponse}
\question{Donner la définition mathématique du courant électrique en précisant les unités.}{1}
\begin{reponse}[18mm]\end{reponse}

\question{Donner la convention du dipôle suivant ainsi que sa relation courant--tension.}{1}
\begin{center}\begin{minipage}{.25\linewidth}\centering\circuitcondensateur\end{minipage}\hfill\begin{minipage}{.68\linewidth}\begin{reponse}[25mm]\end{reponse}\end{minipage}\end{center}

\question{Déterminer l'expression du potentiel $V_A$.}{1}
\begin{center}\begin{minipage}{.35\linewidth}\centering\circuitqfive\end{minipage}\hfill\begin{minipage}{.60\linewidth}\begin{reponse}[36mm]\end{reponse}\end{minipage}\end{center}

\question{Déterminer la résistance équivalente entre $A$ et $B$.}{1}
\begin{center}\begin{minipage}{.38\linewidth}\centering\circuitqsix\end{minipage}\hfill\begin{minipage}{.57\linewidth}\begin{reponse}[28mm]\end{reponse}\end{minipage}\end{center}

\question{Déterminer l'expression de la tension $U_2$.}{1}
\begin{center}\begin{minipage}{.42\linewidth}\centering\circuitqseven\end{minipage}\hfill\begin{minipage}{.53\linewidth}\begin{reponse}[32mm]\end{reponse}\end{minipage}\end{center}

\newpage
\InterroHeader{Outils et lois de l'électrocinétique}{20}
\question{Déterminer l'expression du courant $i_1$.}{1}
\begin{center}\begin{minipage}{.38\linewidth}\centering\circuitqeight\end{minipage}\hfill\begin{minipage}{.57\linewidth}\begin{reponse}[38mm]\end{reponse}\end{minipage}\end{center}

\question{Quelle est la proposition correcte ?}{1}
\begin{tabular}{|p{.47\linewidth}|p{.47\linewidth}|}\hline
A) En convention récepteur, $u(t)=-R\,i(t)$. & D) En convention récepteur, $u(t)=R\,i(t)$.\\\hline
B) En convention générateur, $u(t)=R\,i(t)$. & E) En convention récepteur, $u(t)=R\,i^2(t)$.\\\hline
C) En convention récepteur, $u(t)=i(t)/R$. & \\ \hline
\end{tabular}

\question{La résistance $R_{AB}$ vue aux bornes de $A$ et $B$ vaut :}{1}
\begin{center}
\begin{tikzpicture}[scale=.68]\node[left] at (0,2) {$A$};\node[left] at (0,0) {$B$};\draw (0,2)--(1,2);\draw (1,2)--(1,1.6);\draw (.82,1.6) rectangle (1.18,.4);\node[right] at (1.2,1) {$2R$};\draw (1,.4)--(1,0);\draw (1,2)--(2,2);\draw (2,1.82) rectangle (3,2.18);\node[above] at (2.5,2.2) {$R$};\draw (3,2)--(3,1.6);\draw (2.82,1.6) rectangle (3.18,.4);\node[right] at (3.2,1) {$R$};\draw (3,.4)--(3,0);\draw (0,0)--(3,0);\end{tikzpicture}
\end{center}
\begin{tabular}{|c|c|c|c|c|}\hline A) $R$ & B) $2R$ & C) $4R$ & D) $R/2$ & E) $3R/2$\\\hline\end{tabular}

\question{Pour le diviseur de tension représenté, quelle expression est correcte ?}{1}
\begin{center}\begin{tikzpicture}[scale=.7,>=Latex]\draw (0,0)--(4,0);\draw (0,2)--(.8,2);\draw (.8,1.82) rectangle (2,2.18);\node[above] at (1.4,2.2) {$R_1$};\draw (2,2)--(4,2);\draw (3,2)--(3,1.6);\draw (2.82,1.6) rectangle (3.18,.4);\node[right] at (3.2,1) {$R_2$};\draw (3,.4)--(3,0);\draw[->] (0,.2)--(0,1.8) node[midway,left] {$V_E$};\draw[->] (4,.2)--(4,1.8) node[midway,right] {$V_S$};\end{tikzpicture}\end{center}
\begin{tabular}{|c|c|c|c|c|}\hline A) $V_S=V_E$ & B) $V_S=V_E\frac{R_1}{R_1+R_2}$ & C) $V_S=\frac{V_E}{R_1+R_2}$ & D) $V_S=V_E\frac{R_2}{R_1+R_2}$ & E) $V_S=V_E R_1R_2$\\\hline\end{tabular}

\question{Pour le montage suivant, quelle expression est correcte ?}{1}
\begin{center}\begin{tikzpicture}[scale=.65,>=Latex]\draw (0,0)--(5,0);\draw (0,2)--(.8,2);\draw (.8,1.82) rectangle (2,2.18);\node[above] at (1.4,2.2) {$R$};\draw (2,2)--(5,2);\foreach \x in {2.8,4.6}{\draw (\x,2)--(\x,1.6);\draw (\x-.18,1.6) rectangle (\x+.18,.4);\node[right] at (\x+.2,1) {$R$};\draw (\x,.4)--(\x,0);}\draw[->] (0,.2)--(0,1.8) node[midway,left] {$V_E$};\draw[->] (3.7,.2)--(3.7,1.8) node[midway,right] {$V_S$};\end{tikzpicture}\end{center}
\begin{tabular}{|c|c|c|c|c|}\hline A) $V_S=V_E$ & B) $V_S=V_E/2$ & C) $V_S=V_E/3$ & D) $V_S=V_E/4$ & E) $V_S=2V_E$\\\hline\end{tabular}

\newpage
\InterroHeader{Outils et lois de l'électrocinétique}{20}
\question{Donner l'expression sous forme de somme de sinusoïdes du signal $s(t)$ dont le spectre est donné ci-contre.}{2}
\begin{center}\begin{minipage}{.53\linewidth}\begin{reponse}[52mm]\end{reponse}\end{minipage}\hfill\begin{minipage}{.43\linewidth}\centering\spectrequest\end{minipage}\end{center}

\question{Donner le spectre d'amplitude du signal suivant :
\[i(t)=8+\cos(0{,}02t-0{,}3)+0{,}4\cos(0{,}08t+0{,}5)-0{,}3\cos(0{,}12t-0{,}3).\]}{2}
\begin{reponse}[80mm]\end{reponse}
\end{document}
